\section{引言：为什么需要衡量模型性能}

Yann Dubois 是斯坦福大学三年级博士生，由 Tatsu Hashimoto 和 Percy Liang 共同指导。他在本讲中系统地讲解了 NLP 领域中基准测试（Benchmarking）和评估（Evaluation）的方方面面——这是一个在学术界往往被忽视、但在实际生产中至关重要的话题。

\begin{figure}[H]
\centering
\includegraphics[width=0.95\textwidth]{slides-images/slide-001.jpg}
\caption{课程大纲：涵盖衡量性能的不同原因、封闭式评估、开放式评估（自动评估与人工评估）、当前 LLM 评估方法，以及评估中的问题与挑战\protect\footnotemark}
\end{figure}
\footnotetext{Slides 第2页。}

\subsection{模型开发流水线中的评估需求}

开发一个机器学习模型可以分为四个主要阶段，每个阶段对性能度量有不同的需求：

\begin{figure}[H]
\centering
\includegraphics[width=0.95\textwidth]{slides-images/slide-002.jpg}
\caption{模型开发各阶段对评估的不同要求：从训练到部署，评估的成本容忍度逐渐提高，但对可信度的要求也逐渐增加\protect\footnotemark}
\end{figure}
\footnotetext{Slides 第3页。}

\begin{importantbox}{四个阶段的评估需求}
\begin{enumerate}
    \item \textbf{训练阶段}：需要超快、超廉价、可微分的损失函数，且不能有可被模型利用的``捷径''（shortcuts）
    \item \textbf{开发阶段}（超参调优）：需要快速、廉价的评估指标，同样要避免捷径——因为超参搜索本身也是一种优化
    \item \textbf{模型选择}：可以稍慢、稍贵，但需要在多个候选模型间可靠地比较
    \item \textbf{部署阶段}：评估必须\textbf{可信}（trustworthy）、\textbf{任务特定}（task-specific）、\textbf{绝对}（absolute）——你需要一个明确的阈值来决定是否上线
\end{enumerate}
\end{importantbox}

\begin{knowledgebox}{学术发表 vs 实际部署}
除了上述四个阶段，学术界还有一个``发表''（Publish）需求。学术基准需要：标准化、可复现、易用、低成本、覆盖面广。重要的是，学术基准的指标\textbf{不需要完美}——只要在 10 年的时间尺度上，指标方向与领域实际进步方向一致即可。这与部署场景中对指标精确性的要求截然不同。
\end{knowledgebox}

\subsection{基准测试驱动领域进步}

MMLU（Massive Multitask Language Understanding）是当前最主流的 LLM 基准之一。在过去约 4 年间，模型在 MMLU 上的准确率从 25\%（随机水平，因为是四选一）提升到约 90\%。

\begin{figure}[H]
\centering
\includegraphics[width=0.95\textwidth]{slides-images/slide-003.jpg}
\caption{MMLU 基准上的模型性能进步：从 RoBERTa-base（$\sim$25\%）到 Gemini Ultra（$\sim$90\%），基准测试定义了领域前进的方向\protect\footnotemark}
\end{figure}
\footnotetext{Slides 第4页。}

\begin{warningbox}{学术基准中的细微差异不一定有意义}
在学术基准测试中，重要的不是模型之间 1--2 个百分点的差异，而是在更长时间尺度上的总体趋势。如果过于关注小数点后的改进，可能会导致对噪声的过度拟合。
\end{warningbox}

\subsection{本章小结}

在模型开发的不同阶段，对评估有不同的需求——从训练阶段的快速可微指标到部署阶段的可信绝对指标。NLP 评估可分为封闭式（close-ended）和开放式（open-ended）两大类，每类都有独特的挑战。基准测试是推动领域进步的核心工具，但需要在不同场景下选择合适的评估策略。

%% ========== Part 2: 封闭式评估 ==========

